\documentclass[spanish,xcolor=dvipsnames]{beamer}
    \usefonttheme{professionalfonts}
    \usetheme{Copenhagen}
    \usecolortheme[named=NavyBlue]{structure}
    \setbeamersize{text margin left=1em,text margin right=1.5em}
    \setbeamercovered{transparent}
    \setbeamertemplate{sections/subsections in toc}[circle]
    %\setbeamertemplate{itemize items}[bullet]
    %\setbeamertemplate{itemize subitem}[square]
    %\setbeamertemplate{enumerate items}[square]
    \setbeamertemplate{navigation symbols}{}
    \setbeamertemplate{headline}{}
	
    \usepackage[spanish]{babel}
    \usepackage[T1]{fontenc}
    \usepackage[utf8]{inputenc}
    
	%Estos paquetes pueden ser útiles,
	%si no están instalados puede quitarlos,
	%si los necesita instálelos
    \usepackage{ragged2e}
    \usepackage{lmodern}
    \usepackage[font=scriptsize,labelfont=bf,center]{caption}
    \usepackage{float}
    \usepackage{graphicx}
    \usepackage{subfig}
    \usepackage{multicol}
    \usepackage{amsmath}
    \usepackage{tikz}
    \usetikzlibrary{%
                    automata,%
                    arrows,%
                    shapes,%
                    chains,%
                    matrix,%
                    plotmarks,%
                    %pgfplots.units,%
                    decorations.markings,%
                    decorations.pathmorphing,%
                    backgrounds,%
                    scopes,%
                    calc,%
                    petri,%
                    positioning,%
                    fit%
                    }%
    \usepgflibrary{arrows}
    \usepackage{pgfplots}
    \renewcommand{\shorthandsspanish}{}
    \renewcommand{\spanishtablename}{Tabla}
%    \spanishdecimal{,}

    \makeatletter
        \defbeamertemplate*{footline}{jtheme}
        {
          \leavevmode%
          \hbox{%
          \begin{beamercolorbox}[wd=.25\paperwidth,ht=2.25ex,dp=1ex,center]{author in head/foot}%
            \usebeamerfont{author in head/foot}\insertshortauthor~~\beamer@ifempty{\insertshortinstitute}{}
            {(\insertshortinstitute)}
          \end{beamercolorbox}%
          \begin{beamercolorbox}[wd=.50\paperwidth,ht=2.25ex,dp=1ex,center]{title in head/foot}%
            \usebeamerfont{title in head/foot}\insertshorttitle
          \end{beamercolorbox}%
          \begin{beamercolorbox}[wd=.25\paperwidth,ht=2.25ex,dp=1ex,right]{date in head/foot}%
            \usebeamerfont{date in head/foot}\insertshortdate{}\hspace*{2em}
            \insertframenumber{} / \inserttotalframenumber\hspace*{2ex}
          \end{beamercolorbox}}%
          \vskip0pt%
        }
    \makeatother


    %%-------------------Portada-------------------------%%
    \title{Título de la presentación}
    \author[Nombre Corto]{Nombre Completo}
    \institute[ULA]{
        Facultad de Nombre\\
        Departamento o Escuela
    }
    \date[Mes Año]{\footnotesize Mérida, Mes, Año}

    %%-------------------Fin-Portada---------------------%%

\begin{document}
	%%--------------------Diapositiva--------------------%%
    \begin{frame}[plain] % [plain] para que no salga el contenido en esta página
	    \begin{center}
	    	\includegraphics[width=1.5cm]{ULA_logo_titulo}
	    \end{center}
	    \titlepage
    \end{frame}
    %%------------------------Fin------------------------%%
	
	%Logo de la esquina superior derecha de cada frame, necesita tres compilaciones para actualizar.
	\addtobeamertemplate{frametitle}{}{%
    \begin{tikzpicture}[remember picture,overlay]
    	\node[anchor=north east,yshift=11pt,xshift=11pt] at (current page.north east) {\includegraphics[height=1.5cm,angle=-20,keepaspectratio]{ula}};
    \end{tikzpicture}}
    
	%%--------------------Diapositiva--------------------%%
    \begin{frame}\justifying
        \frametitle{Agenda}
        \begin{columns}
            \begin{column}{0.0cm}
            \end{column}
            \begin{column}{0.9\textwidth}
                \tableofcontents
            \end{column}
        \end{columns}
    \end{frame}
    %%------------------------Fin------------------------%%

    %%--------------------Diapositiva--------------------%%
    \section{Primera sección}
    \begin{frame}
        \frametitle{Título del frame (no obligatorio)}
        \framesubtitle{Subtítulo del frame (no obligatorio)}
        \begin{columns}
        	\begin{column}{.7\linewidth}
        		Contenido del frame a dos columnas de distinto ancho.
        	\end{column}
        	\begin{column}{.3\linewidth}
        		Contenido del frame a dos columnas de distinto ancho.
        	\end{column}
        \end{columns}
    \end{frame}
    %%------------------------Fin------------------------%%

    %%--------------------Diapositiva--------------------%%
    \begin{frame}
        \frametitle{Título del frame}
        \framesubtitle{subtítulo del frame}
        \begin{theorem}
        	Algún teorema
        	
        	\begin{equation}
        		f(x) = posible\qquad ecuación \qquad teorema
        	\end{equation}
        	
        \end{theorem}
    \end{frame}
    %%------------------------Fin------------------------%%
    
	%%--------------------Diapositiva--------------------%%
    \begin{frame}
        \frametitle{Título del frame (no obligatorio)}
        \framesubtitle{Subtítulo del frame (no obligatorio)}
        \begin{block}{Ejemplo de bloque}
        	Contenido del bloque. \textcolor{green}{Texto de otro color}.
        \end{block}
    \end{frame}
    %%------------------------Fin------------------------%%
    
	%%%%%%%% % % % % % % % % % % % % % % % % % % %%%%%%%%%%
	%%%%				Use para nuevo frame		   %%%%
	%%%%%%%% % % % % % % % % % % % % % % % % % % %%%%%%%%%%
   	%%--------------------Diapositiva--------------------%%
    %\begin{frame}
    %    \frametitle{¿Mínimos cuadrados?}
    %    \framesubtitle{Método de optimización matemática}
        
    %\end{frame}
    %%------------------------Fin------------------------%%
    
    %%--------------------Diapositiva--------------------%%
    \section{Bibliografía}
    \begin{frame}
        \frametitle{Ejemplo de Bibliografía con item alfanumérico}
        \begin{thebibliography}{9}\small
            \setbeamertemplate{bibliography item}[text]
            \bibitem{A} Proakis, John y Malonakis, Dimitris. "Least-Square Design Methods". Capítulo: \emph{Design of Digital Filters}. \emph{Digital Signal Processing. Principles, Algorithms and Applications}. Prentice Hall International. 1996. pp. 706--714.
            \bibitem{B} "Least-Squares Fitting - MATLAB \& Simulink". \emph{Mathworks Documentation Center}. Consulta: 09/06/2014. URL: \url{http://www.mathworks.com/help/curvefit/least-squares-fitting.html}.
            \bibitem{C} "Mínimos Cuadrados". \emph{Wikipedia, La Enciclopedia Libre}. Consulta: 09/06/2014. URL: \url{http://es.wikipedia.org/wiki/M\%C3\%ADnimos_cuadrados}.
            \bibitem{D} "Recursive digital filters design - MATLAB yulewalk". \emph{Mathworks Documentation Center}. Consulta: 09/06/2014. URL: \url{http://www.mathworks.com/help/signal/ref/yulewalk.html}.
        \end{thebibliography}
    \end{frame}
    %%------------------------Fin------------------------%%
\end{document}